\section{Method}

\paragraph{Overview.}
\textsc{Heavy-Hitter Masking} is implemented as a Hugging Face data collator for masked language modeling. It receives tokenized examples, selects valid non-special positions to mask, applies the same replacement policy as the random baseline, and returns metadata used to update the online error sketch. The method is intentionally placed in the collator rather than in the model: this makes the intervention cheap, easy to disable, and directly comparable to random masking. The collator is also the right level of abstraction for the scientific question. Prior adaptive masking work often improves a full training recipe by combining mask selection with schedules, auxiliary regularization, or representation changes \citep{yang2023learning,zhong2023self,edman2025mask}; here the collator boundary enforces that the model and optimizer are unchanged.

\paragraph{Masking Mixture.}
HHM masks the same total proportion of tokens as the baseline: 15\% of valid positions. The default mixture allocates 50\% of masked positions to uniform random masking, 25\% to local entity heavy hitters, and 25\% to online error heavy hitters. The 50\% random component is a guardrail: it keeps broad coverage over the corpus and prevents the adaptive policy from collapsing onto a narrow set of tokens. If either heavy-hitter source lacks enough positions, the remainder is filled by random valid positions. This fallback is part of the method rather than an implementation accident, because some batches simply do not contain enough repeated or previously hard tokens to support the target mixture.

\paragraph{Local Entity Heavy Hitters.}
For each sequence, a lightweight heuristic extracts entity-like token strings: alphabetic non-stopword tokens that are at least two characters long and either capitalized or content-word-like. A Space-Saving sketch with capacity 64 counts candidate strings. Positions belonging to repeated top entities are eligible for the entity-guided mask share. This component was tested because BabyLM evaluation includes entity-tracking diagnostics, and because repeated mentions are a place where a small model may benefit from seeing the same discourse participant under multiple contexts. The heuristic is deliberately shallow: the goal is not to solve entity recognition, but to ask whether a cheap recurrence signal is useful at all.

\paragraph{Online Error Heavy Hitters.}
After warmup, the training loop computes per-token MLM losses for masked positions. The top-loss masked tokens in each batch update a bounded Space-Saving sketch. Later batches preferentially mask positions whose token string appears in this error sketch. This tests whether the model can provide its own training signal: instead of assuming in advance which tokens matter, HHM lets recent prediction failures influence future masking. Warmup prevents the sketch from being dominated by the model's earliest, largely uninformative errors. The negative HHM and CMS-Morph results show why this raw-loss signal is risky: it can overemphasize rare, ambiguous, or tokenization-sensitive items rather than stable learnable regularities.

\paragraph{CMS-Morph Pilot.}
CMS-Morph is an additional pilot variant of the same masking-only idea. It maintains two Count-Min Sketches: one over high-loss token strings and one over character trigram features extracted from high-loss tokens. The motivation is that a token-only sketch may memorize isolated hard words without transferring that signal to related forms. Character trigrams provide a small amount of subtoken generalization: if forms sharing a suffix or stem-like pattern are repeatedly hard, future examples with the same pattern can receive part of the mask budget. During training CMS-Morph keeps the total mask rate fixed at 15\%, but after warmup allocates 70\% of masked positions uniformly at random, 20\% to token-sketch-guided positions, and 10\% to trigram-sketch-guided positions. Sketches are updated only from masked-token losses in the training batches; validation and test data are never used for sketch updates.

\paragraph{Accuracy-Morph Pilot.}
Accuracy-Morph is a second pilot motivated by a failure mode of raw loss: high loss can reflect noise, ambiguity, or rare-token scale rather than a stable learnable difficulty. Accuracy-Morph therefore replaces raw high-loss updates with smoothed prediction-correctness statistics. For each masked token, the training loop records whether the model's top prediction matches the MLM label. Two token sketches count masked-token exposures and wrong predictions, and two analogous sketches count character trigram exposures and wrong predictions. Candidate positions are scored by smoothed error rate,
\[
  \frac{\mathrm{wrong} + 0.5}{\mathrm{seen} + 1.0},
\]
rather than raw loss, so a token or trigram becomes attractive only when it has been observed enough times and is often predicted incorrectly. The pilot uses a longer 30\% warmup and a conservative 80/15/5 random/token/trigram mask mixture, again preserving the 15\% total mask rate. This variant was tested to separate the idea of adaptive masking from the noisiness of raw-loss heavy hitters.

\paragraph{Entity-Accuracy-Morph Pilot.}
Entity-Accuracy-Morph tests whether the promising Accuracy-Morph signal can be improved by borrowing the repeated-token intuition from the strongest prior BabyLM systems. It uses the same correctness sketches and the same 30\% warmup as Accuracy-Morph, but adds a small score bonus for repeated non-special token strings within a batch and restricts adaptive candidates to an intermediate predicted-error band. The band restriction was intended to avoid two extremes: tokens that are already too easy to be useful and tokens that may be too noisy or idiosyncratic to learn from. The adaptive share is also decayed late in training, moving from an 80/15/5 random/token/trigram mixture toward a 90\% random mixture after 80\% of training. This variant is intended to test whether entity-like repetition can help the correctness-smoothed masking signal without changing the model or total mask rate.

\paragraph{Design Principle.}
The method deliberately avoids changing the model, tokenizer, dataset, optimizer, schedule, or total mask rate. This keeps the comparison focused on the masking policy.
